% !TeX program = xelatex
% Overleaf: Menu → Compiler → XeLaTeX  （或使用本目录 latexmkrc）
%
% 换期刊模板时主要改这一段 documentclass / 导言宏包，正文 \cite / \section 可尽量保留。
% Elsevier 例：
%   \documentclass[preprint,12pt]{elsarticle}
% IEEE 例：
%   \documentclass[journal]{IEEEtran}
%
% ctexart 的默认中文字体集 fandol 只支持 XeLaTeX / LuaLaTeX。
% 下面这行让"用错编译器"时直接给出可读的提示，而不是 ctex 的 fontset 报错。
\RequirePackage{iftex}
\RequireXeTeX
\documentclass[UTF8,a4paper,11pt]{ctexart}

\usepackage[margin=2.2cm]{geometry}
\usepackage{amsmath,amssymb,amsfonts,mathtools,bm,amsbsy,amsthm}
\usepackage{graphicx}
\usepackage{pifont}
\usepackage{booktabs,tabularx,array,makecell,multirow}
\usepackage{algorithm}
\usepackage[section]{placeins}
\usepackage{xcolor}
\usepackage{setspace}
\usepackage[backend=biber,style=numeric,sorting=none,maxbibnames=99]{biblatex}
\addbibresource{refs.bib}
\usepackage{hyperref}
\usepackage[nameinlink]{cleveref}

\hypersetup{colorlinks=true,linkcolor=blue,citecolor=blue,urlcolor=blue}
\graphicspath{{figures/}{./}}
\allowdisplaybreaks
\newtheorem{assumption}{假设}
\newtheorem{remark}{注}
\newtheorem{problem}{问题}
\DeclareMathOperator{\diag}{diag}

% MathType AMS-LaTeX 常见残留：多余花括号、array 列格式 *{20}{l}
